\documentclass[10pt]{article}
\usepackage[margin=0.8in]{geometry}
\usepackage{amsmath,amssymb,booktabs,longtable,hyperref}
\usepackage[T1]{fontenc}
\hypersetup{colorlinks=true,urlcolor=blue,linkcolor=blue}
\title{Flow-Initialized Bounded RTI with One Analytic Quadratic CLF}
\author{Collision Avoidance Research}
\date{October 2, 2026}
\begin{document}
\maketitle

\section{Outcome and measurement scope}
The final controller avoided collision and recovered given-path constant-speed
cruise in all fourteen deterministic threat scenarios. Each uncontrolled,
given-path cruise baseline contains a collision. Success requires actual
recovery, rather than survival for a fixed short duration. The campaign executed
5755 holds: median full controller-call time 12.788 ms, 95th percentile 25.526 ms,
and maximum continuation-call time 87.949 ms. The smallest independently refined
physical rectangle clearance was 0.0787976 m.

Of these calls, 5670 (98.52\%) built one trajectory model and 5604 (97.38\%) used
one numerical solve. The remaining 85 calls built two models; none built three.
Initialization and restoration can require extra solves: the largest count was
six, despite the two-model bound. Twenty calls accepted damping, using 43 cheap
trial rollouts in total. Every issued result completed CLF optimization and
satisfied the existing prediction-agreement tolerances.

These timings are measured computation times, not just native solver times.
The ODE45 plant replay does not apply the measured computation delay to the
plant. The nominal sample interval remains 50 ms, and 35 calls exceeded 50 ms.
The two calls exceeding 100 ms were the initial 8 and 15 m/s head-on calls:
231.377 and 122.200 ms. A separate fresh MATLAB process measured a 1409.726 ms
first call, including loading, JIT and terminal preparation. The six subsequent
fresh-problem calls in that process had median 37.715 ms. Initialization is
therefore not within 100 ms, and the measured campaign is not a hard deadline
certificate. No deadline, state constraint, collision buffer or solver tolerance
was relaxed to obtain these timings.

\section{Implemented controller}
A shifted input sequence is rolled out from the measured state; an unusable or
missing sequence uses the moving-target flow initialization. The reference is a
nonlinear Fiala bicycle trajectory. Each trajectory model shares the same anchor
for dynamics, tire derivatives, collision rows, terminal geometry and CLF.
The ordinary path is one model and one CLF solve under the inherited PCBF budget:
\[
 \sum_i \xi_{i|k}\leq\sum_{i=1}^{N-1}\xi_{i|k-1},\qquad
 \xi_{0|k}\leq\xi_{1|k-1}.
\]
Primary minimization remains available to initialize or restore that problem.
It does not supply an executable primary-only or previous-plan backup.

\paragraph{One analytic CLF.}
The only active nominal function is
\[
 V(x)=e(x)^\top P e(x),\quad
 e=[e_y,e_\psi,v_x-v_{x,\rm ref},v_y-v_{y,\rm ref},r-r_{\rm ref}]^\top.
\]
P is the discrete LQR matrix of the transverse cruise linearization. The desired
inequality is $V(x_1)-V(x_0)\leq-\eta e_0^\top Qe_0+\rho$, with $\eta=0.5$.
The analytic error Jacobian on the affine successor gives
$V_{\rm aff}=\|\operatorname{chol}(P)(e(\bar x_1)+J_e\Delta x_1)\|^2$.
This has positive semidefinite Hessian and one SOC epigraph. The terminal cone
remains, so the problem is a sparse conic QP/SOCP, not a QP with only affine
constraints. The expensive nominal return rollout, finite-difference Hessian,
old cost-to-go utility, its compiled-kernel build and obsolete configuration
fields have been removed. Path guidance supplies only initialization inputs.

\paragraph{Feasible damping.}
Before a new model is built, the complete convex decision vector is interpolated
between a verified feasible primary/inherited point and the CLF optimizer point.
The inherited zero correction is not assumed feasible merely because it has a
slack budget. The base must satisfy the exact assembled equalities, inequalities,
boxes and cones. The CLF epigraph is tightened by evaluating the same quadratic
at the trial point. At most three fractions are tried, starting at
$\min(0.5,0.8/\sqrt{r})$ and halving. Each trial uses a nonlinear rollout, without
a new solve or new derivatives. A zero-slack full CLF solution must remain zero
within tolerance after damping. With positive optimal slack, damping can sacrifice
some improvement; the issued slack and applied fraction are reported explicitly.

\paragraph{Two actual models at most.}
If damping is unavailable or inaccurate, one additional model may be built.
The new correction radius is at most half the old radius. A fresh-flow retry
shares this same allowance. Restoring a failed inherited budget reuses the
existing matrices and CLF cone. A bounded input-box enlargement changes only
bounds and the normalized input-increment objective, not the linearization.
This fixes an otherwise hidden repeated model build. Configuration validation
rejects more than two builds. Inputs are still penalized about the anchor rather
than zero, and steering has no physical magnitude or slew constraint.

\section{Why the intermediate designs were insufficient}
Adding damping to the old cost-to-go CLF passed all fourteen cases, but only 37
of 6532 calls accepted a damped step. Its median was 29.396 ms and its worst
noninitial recorded call was 189.523 ms. Damping alone did not remove CLF
construction cost or make the normal architecture bounded to two models.

A first quadratic-CLF prototype with two refinement rounds failed six 15 m/s
cases. Replaying those failed states with the diagnostic eight-round allowance
succeeded in exactly three rounds. The inherited interpolation bases violated
only the tiny terminal cone, so damping could not legally start from them.
Ratios just above one induced only modest radius reductions, sometimes increasing
the next model error. A half-radius correction and reuse of matrices resolved
these cases without adding another controller or permitting a third model.
A later strict model-count prototype exposed the low-speed braking initialization:
input-box enlargement had consumed the second model build despite an unchanged
anchor. Reusing those matrices restored the final fourteen-case result.

The practical lesson is that a feasible interpolation segment is conditional.
A flow or shifted trajectory is a search anchor, not automatically a feasible
point of the current terminal-constrained problem. The implementation does not
hide that condition or treat damping as a universal nonlinear safety guarantee.

\section{Closed-loop results}
Exact states are used, with no observer error and no road-boundary constraints.
The target keeps constant tangential acceleration and sideslip. Speeds are 8 and
15 m/s; PCBF prefixes have eight and sixteen holds, followed by a three-second
completion tail. The affine collision-node buffer remains 0.10 m. Independent
ODE45 rectangle replay measures physical clearance; positive clearance is the
collision-free criterion. Recovery means all five errors are within
$[0.1\,\mathrm m,1^\circ,0.1\,\mathrm{m/s},0.05\,\mathrm{m/s},0.01\,\mathrm{rad/s}]$
for one second after the baseline conflict. The 1500-hold observation cap is
not a success deadline. There is no random seed because these fixtures are deterministic.

\begin{longtable}{rlrrr}
\toprule
Speed & Scenario & Min. gap (m) & Recovery (s) & Max. continuation (ms)\\
\midrule\endhead
8 & headOn & 1.9247 & 20.75 & 60.29 \\
8 & acceleratingHeadOn & 1.5191 & 21.55 & 57.30 \\
8 & brakingLead & 0.0995 & 14.55 & 87.95 \\
8 & crossing & 0.6592 & 16.65 & 46.68 \\
8 & turningCrossing & 0.7260 & 17.90 & 48.80 \\
8 & curvedHeadOn & 2.3730 & 24.65 & 46.62 \\
8 & curvedCrossing & 0.4353 & 20.30 & 47.67 \\
15 & headOn & 0.2144 & 11.60 & 41.37 \\
15 & acceleratingHeadOn & 0.0788 & 11.20 & 46.40 \\
15 & brakingLead & 0.0993 & 15.90 & 79.28 \\
15 & crossing & 2.6877 & 13.05 & 45.46 \\
15 & turningCrossing & 2.6876 & 13.05 & 40.25 \\
15 & curvedHeadOn & 2.0597 & 70.10 & 59.90 \\
15 & curvedCrossing & 1.8252 & 16.50 & 46.30 \\
\bottomrule
\end{longtable}
The listed gaps use the standard replay mesh. The three smallest-gap cases were
independently reintegrated over their closest hold and both neighbors, with 1001
samples per hold (50 microseconds), ODE45 relative tolerance $10^{-12}$ and
absolute tolerance $10^{-13}$. Their refined minima were 0.0787976 m for the
15 m/s accelerating head-on case, 0.0992696 m for the 15 m/s braking case, and
0.0995284 m for the 8 m/s braking case. These are numerical replay observations,
not interval-certified physical lower bounds.

The independent Python geometry calculation agrees with MATLAB in all fourteen
cases. All 5668 inherited-budget calls satisfy the transfer equations; the largest
sum-transfer discrepancy is $4.24\times10^{-22}$. Maximum optimized budget excess
is $7.55\times10^{-10}$. All selected numerical stages terminated positively;
885 stages used the solver's declared reduced-accuracy termination. Twelve calls
have PCBF slack above $10^{-5}$; the largest total is $1.634\times10^{-4}$.
They remain relaxed predictions, even though physical replay was collision-free.
The legacy audit's combined ``passed'' flag uses a different hard-feasibility
contract; only its independent geometry/recovery outputs are used here.

\section{Runtime composition and paired comparison}
The slowest continuation was the 8 m/s braking scenario at $t=5.35$ s. It built
two models and made five numerical solves. Disjoint measured components are:
\begin{center}\begin{tabular}{lr}\toprule
Component & Milliseconds\\\midrule
Reference rollout & 0.609 \\
Dynamics, collision and terminal formulation & 20.459 \\
Analytic CLF construction & 0.277 \\
Primary PCBF solves & 15.412 \\
CLF solves & 44.514 \\
Full-step nonlinear agreement & 1.284 \\
Feasible-segment line search & 2.084 \\
Other orchestration and output & 3.310 \\
\midrule Total & 87.949\\\bottomrule
\end{tabular}\end{center}
Native PCBF and CLF stages occupy 59.926 ms (68.1\%); formulation excluding CLF
occupies 20.459 ms (23.3\%). Analytic CLF construction is only 0.277 ms.
This remaining exceptional cost is mostly solver work and two trajectory models,
not repeated nominal-value simulation.

A same-process comparison used six alternating pairs after three warmups, with
the terminal object prepared outside each timed call. The original 8 m/s head-on
initial state needed two rounds and three solves: median 79.157 ms. The quadratic
RTI version needed one model and one solve: 27.823 ms. At the original slow braking
state ($t=5.45$ s), the exact historical controller/continuation was reconstructed
with input error below $10^{-9}$; medians were 254.658 ms (four rounds, nine solves)
and 28.793 ms (one model, one solve). The state and prior plan are identical in
each pair. Closed-loop future states need not remain identical after issuing
the changed control.

Those pairs were recorded immediately before deleting inactive cost-to-go
utilities and their configuration fields; the active quadratic, damping,
constraints and solver were already the final algorithms. Frozen measured
sources are preserved with the diagnostic export. Campaign comparisons across
processes are not interpreted as isolated causal timing measurements. External
system load, JIT and cache state affect wall time.

\section{Mathematical limits and validation}
The LQR identity supplies strict decrease for the local linear system; a smooth
nonlinear local result also needs an admissible neighborhood of the sampling
equilibrium. It does not give global decrease for large lateral/heading errors.
The former cost-to-go reachable-domain argument is not claimed for this quadratic.
At a 20-m offset, four old zero-slack assertions fail: positive slack is needed
although V decreases. They were replaced by tests of the actual stated scope:
local zero-slack decrease and eventual recovery from displaced states.
Eight closed-loop starts cover 20-m displacement and 40-m displacement with a
2.5-radian heading error, both speeds and both curvatures; all return to cruise
and hold the tolerances. No fixed recovery speed is asserted.

There are 393 passing related MATLAB unit tests across eighteen classes, plus
120 passing signed local-quadratic decrease samples. Tests cover model and
geometry kernels, shifted budgets, restored budgets, damping, actual two-model
limits, no inaccurate-command issuance, same CLF across target ranges, terminal
membership, solver termination and nonlinear recovery. The model/collision
agreement threshold is not a complete nonlinear constraint certificate, and the
encounter-range contract is unchanged. A retained positive-slack prediction is
not a proof of nonlinear recursive feasibility. Four damping attempts correctly
refused to give up a full-step zero CLF slack.

\section{Artifacts and reproduction}
Current mathematical definitions are in \path{controller/NOMINAL_CLF.md} and
\path{controller/PCBF_CLF_ARCHITECTURE.md}. The accompanying
\path{report/FLOW_RTI_QUADRATIC_20261002/} directory contains compact scenario
results, timing records, the worst continuation trace, independent audits,
local decrease samples and source hashes. Full traces, intermediate failed
experiments, executed analysis scripts and frozen paired sources remain in the
identified cache and external diagnostic export; they are not vendored into the
controller or Git repository.

From the repository root, use MATLAB R2026a Update 3 with
\texttt{-singleCompThread}; add \texttt{scripts}, \texttt{controller},
\texttt{config} and \texttt{solver/controller}. Run
\texttt{runTimedFlowCampaign(pwd,outputDirectory,1500)} for the fourteen scenarios.
Run \texttt{validateNominalClf} for the local decrease grid. Related unit-test
class lists and actual commands are preserved in the activity record and test
logs. The optional compiled model/geometry kernels were unchanged and tested;
\texttt{buildControllerKernels} no longer builds a nominal cost-to-go kernel.
The final Code Analyzer reports eight sparse-indexing performance advisories in
formulation rows and no messages in the other changed MATLAB files. Whitespace
validation and standalone report compilation pass. No generated binary or
unrelated estimator/manuscript change is included.
\end{document}
